\section{Hubs and Targeted Damage}
\label{sec:hubs}

\subsection{Where the extra loss comes from}
\label{sec:hubs-loss}

Removing the 5\% of nodes with the highest degree from the full graph leaves
80.4\% of nodes in the giant component, against 93.8\% after random removal of
the same number of nodes (Figure~\ref{fig:summary}c). The indications lose 82.8\% of
their two-hop links, against 12.3\% for random removal, as Eq.~\eqref{eq:r2-node}
gives (12.4\%). The giant component registers that the attack is worse, but
it understates the difference by a factor of five.

The loss has two parts of similar size. Of the 82.8 points, 43.6 come from lost
endpoints, almost all of them diseases, and 39.2 from lost routes of
indications whose drug and disease both survive. Lost endpoints are a selection
effect. Having a shared protein favors diseases with many protein links: the
median two-hop indication has a disease of degree 88, against 22 for all
indications, and the 5\% attack removes nodes down to degree 108.

\subsection{Size-biased routes and the mean field}
\label{sec:hubs-mf}

Table~\ref{tab:mf} tests the mean field of Eq.~\eqref{eq:mf} and separates its
two assumptions. The first column uses the graph-wide route share $h$; the
second replaces it by the share of the task's own routes that the attack cuts,
which removes the selection error and leaves only the correlation error.

\begin{table}[!htbp]
\centering
\small
\caption{Route survival under degree attack for queries whose endpoints survive.
``Graph-wide'' uses Eq.~\eqref{eq:mf} with $h$ computed from all
drug--protein--disease routes of the graph; ``own routes'' uses the share of
the task's own routes that the attack cuts. Full graph: the 1,931 two-hop
indications. Sample: the 200 benchmark queries.}
\label{tab:mf}
\resizebox{\linewidth}{!}{%
\begin{tabular}{@{}llccccc@{}}
\toprule
& Removed & $h$ (graph) & $h$ (own routes) & Mean field (graph) & Mean field (own routes) & Measured \\
\midrule
Full graph & 0.5\% & 0.31 & 0.21 & 0.780 & 0.858 & 0.853 \\
           & 1\%   & 0.44 & 0.27 & 0.654 & 0.798 & 0.784 \\
           & 2\%   & 0.60 & 0.45 & 0.497 & 0.643 & 0.642 \\
           & 5\%   & 0.78 & 0.74 & 0.283 & 0.338 & 0.305 \\
           & 10\%  & 0.89 & 0.81 & 0.134 & 0.227 & 0.212 \\
\addlinespace
Sample     & 1\%   & 0.39 & 0.37 & 0.764 & 0.784 & 0.632 \\
           & 2\%   & 0.51 & 0.42 & 0.665 & 0.742 & 0.587 \\
           & 5\%   & 0.67 & 0.67 & 0.477 & 0.484 & 0.462 \\
           & 10\%  & 0.77 & 0.75 & 0.351 & 0.383 & 0.370 \\
\bottomrule
\end{tabular}}
\end{table}

The two graphs fail the mean field in opposite ways. On the full graph, the
routes of real indications are cut almost independently: with their own route
share, the mean field is within 0.034 at every attack size. The graph-wide
mean field fails because of selection, and direct counting shows what is
selected: at a 1\% attack, 27\% of the indications' routes are cut, against 44\%
of all routes in the graph and 51\% of the routes of two-hop decoys, which have
the same degree distributions. Real indications avoid hub proteins; chance
pairs seek them out. The mean field therefore underestimates the survival of
indications by up to 0.15, and its agreement at 5\% is a cancellation of errors. On the sample, the routes of the benchmark are typical
($h$ is nearly the same for both columns at 5\% and 10\%), but small attacks
cut the routes of a query together, and the mean field overestimates survival by
about 0.15 at 1\% and 2\%. In the sample, low-degree nodes hang from hubs
(Section~\ref{sec:data-sample}), so a drug that reaches its disease through one
hub tends to do so only through hubs. The size bias itself is plain in both graphs. On the full graph the 1\% of nodes
with the highest degree carry 44\% of all drug--protein--disease routes, and in
the sample the median route protein has degree 188, against 25 for all
proteins.

\subsection{Hubs carry the non-specific paths}
\label{sec:hubs-specific}

Hub proteins are where chance paths concentrate. Of the shared proteins that
link decoys, 45\% are among the 1\% of proteins with the highest degree (degree
at least 316); for indications the share is 24\%. Removing proteins alone, in
order of degree, isolates this effect, because no drug or disease is removed
(Figure~\ref{fig:indications-attack}b). Removing the top 155, 310 and 619
proteins lowers two-hop recall from 31.0\% to 26.4\%, 23.7\% and 21.3\%, but it
raises the ratio of indication to decoy rates from 2.18 to 2.96, 3.33 and 3.87,
and the information per pair from 29.6 to 39.2, 40.0 and 41.0 millibits. At
$L=3$ the information rises from 11.9 to 16.4 millibits. Removing 619 proteins
at random instead leaves the ratio at 2.17, so the gain is specific to hubs. A
softer version, scoring each pair by the sum of $1/k$ over its shared proteins,
keeps all the recall but barely improves the separation (AUROC 0.591 against
0.586), so the specificity gained by removing hubs is paid for in recall. Hubs
therefore supply
much of the redundancy that keeps short paths alive under blind damage, and
many of the paths that connect unrelated pairs. Network medicine corrects for
this by comparing a drug--disease proximity with degree-preserving random
references \cite{guney2016network}, and link prediction in protein networks
discounts paths through high-degree nodes \cite{kovacs2019network}.

\begin{figure}[!htbp]
  \centering
  \includegraphics[width=0.86\linewidth]{act6/fig_indications_attack.pdf}
  \caption{Degree attack on the full graph. (a) Giant component, and recall of
  indications (solid) and decoys (dotted) at $L=2,3,4$, when nodes of any type
  are removed in order of degree. (b) Ratio of indication recall to decoy rate,
  on a logarithmic axis, for the attack on all nodes (solid, up to 10\% of
  nodes) and on proteins only (dashed, $L=2$ and $L=3$, up to 5\%).}
  \label{fig:indications-attack}
\end{figure}
